% Заключение.

\thesisChapterStar{Заключение}

В настоящей выпускной квалификационной работе спроектирован и реализован комплекс клиентских приложений SaaS-платформы <<ВКурсе>> для управления учебным расписанием, обеспечивающий единый пользовательский опыт работы с расписанием для студентов, преподавателей и сотрудников учебных заведений. Поставленная цель достигнута; в ходе работы решены пять сформулированных во введении задач.

\textbf{Анализ предметной области и формирование требований.} Проведён обзор российского рынка систем управления учебным расписанием с разделением на два класса: прямые конкуренты (<<Галактика РУЗ>>, <<ММИС~Лаб>>, <<1С:Университет>>), ориентированные на учебное заведение как заказчика и поставляемые по on-premise-модели, и непрямые конкуренты --- мобильные приложения-парсеры, агрегирующие расписание из публичных источников. На стороне ручного ведения расписания выделены три класса системных потерь --- ресурсные коллизии при составлении, задержки в доставке оперативных изменений и отсутствие обратной связи об использовании. По итогам анализа выделены три целевые роли (студент, преподаватель, методист), сформулированы одиннадцать функциональных и восемь нефункциональных требований к клиентской части платформы.

\textbf{Проектирование клиентской части.} В основу проектирования положен принцип максимальной нативности: каждое клиентское приложение следует рекомендациям своей целевой платформы (Apple Human Interface Guidelines, Material~Design~3 Expressive, shadcn/ui поверх Tailwind~CSS) с собственной дизайн-системой и отдельным Figma-файлом. Спроектирован комплекс из четырёх клиентских компонентов, объединённых единым серверным контуром через принцип единого фасада: iOS- и Android-приложения, веб-приложение просмотра расписания с лендингом и закрытый веб-сервис администрирования. На уровне архитектуры клиентов введён кросс-платформенный изоморфизм слоёв при сохранении нативной идиоматики каждой платформы; для редактора расписания спроектированы атомарная модель урока и двухфазная транзакционная схема публикации изменений с серверной валидацией ресурсных коллизий.

\textbf{Реализация мобильных приложений.} iOS-приложение реализовано на Swift и SwiftUI с использованием фреймворка Observation для реактивности, SwiftData --- для локального хранения и WidgetKit --- для виджетов рабочего стола; релиз выпущен с поддержкой 12 акцентных цветов, 8 вариантов иконки приложения и двойной визуальной адаптации под визуальный язык Liquid~Glass на iOS~26 и выше. Android-приложение реализовано на Kotlin и Jetpack Compose с применением Material Design~3 Expressive, Room в качестве реляционного хранилища, StateFlow для реактивных моделей и Navigation Compose с типобезопасными маршрутами; поддерживаются 9 акцентных схем и 8 иконок приложения через механизм \texttt{activity-\allowbreak alias}. Оба клиента построены вокруг общего композиционного корня из восьми доменных моделей-наблюдателей, выполняют глубокое упреждающее кэширование расписания в окне $\pm 1$~год и обеспечивают полноценную офлайн-работу. Системные виджеты рабочего стола реализованы на iOS; их реализация на Android вынесена в задачи дальнейшего развития.

\textbf{Реализация веб-приложения.} Веб-составляющая реализована как единый Next.js-проект (App Router) с тремя функциональными ролями, разделёнными маршрутными группами: маркетинговый лендинг с MDX-контентом для блога и сопутствующих страниц, публичное веб-приложение просмотра расписания с mobile-first-вёрсткой и параметризацией через URL для шаринга ссылок на расписание сущности, а также закрытый веб-сервис администрирования с авторизацией по схеме OAuth~2.0 в режиме ROPC и парой JWT-токенов. Ключевая архитектурная особенность веб-составляющей --- собственный BFF-слой средствами обработчиков маршрутов Next.js, через который проходят все обращения к серверной части и который скрывает общий клиентский ключ от браузера, реализуя требование к безопасной обработке токенов. В сервисе администрирования реализованы управление справочными сущностями учебного заведения через TanStack~Table и редактор расписания на основе атомарной модели урока со стеком undo/redo и двухфазной публикацией изменений.

\textbf{Апробация результатов.} Все три клиентских компонента опубликованы в крупнейших каналах распространения: iOS-приложение --- в App~Store (сентябрь 2025~г.), Android-приложение --- в Google~Play и RuStore (февраль 2026~г.), веб-составляющая --- на собственном домене \texttt{vcourse.app} (февраль 2026~г.). Суммарное количество установок мобильных приложений превышает 3000, ежедневная активная аудитория составляет порядка 1000 уникальных пользователей, распределение по платформам --- iOS~79\%, Android~19\%, Web~1\%. К парсинговому контуру серверной части подключены шесть омских учебных заведений; наибольшая нагрузка приходится на Омский государственный университет (около 4840 обращений к API в сутки) и Сибирский государственный автомобильно-дорожный университет (около 2562 обращений). Сопровождение проекта в публичном поле ведётся через Telegram-канал <<ВКурсе: Расписание занятий>> с аудиторией более 410 подписчиков. Проект представлен инвестиционному фонду <<ТРАМПЛИН>> (Омск, апрель 2026~г.) и во втором этапе грантового конкурса <<Студенческий стартап>> Фонда содействия инновациям (май 2026~г.). 15~мая 2026~года в Омском государственном аграрном университете состоялась первая очная встреча с потенциальным пилотным партнёром; согласована повторная встреча на середину июня 2026~года и переход к обсуждению условий пилотного внедрения.

\textbf{Научная новизна и практическая значимость.} Научная новизна работы заключается в комплексном подходе к построению клиентской части SaaS-платформы для управления учебным расписанием, отличающемся от существующих российских решений одновременным сочетанием SaaS-модели поставки, нативной мобильной составляющей и независимости от сторонних ERP-экосистем, а также в архитектурном решении --- кросс-платформенном изоморфизме слоёв клиентских приложений при сохранении нативной идиоматики каждой целевой платформы. Практическая значимость работы подтверждается фактической публикацией клиентских приложений в трёх крупнейших каналах распространения, устойчивой ежедневной активной аудиторией порядка 1000 уникальных пользователей, обработкой реальной пользовательской нагрузки шести подключённых учебных заведений и началом институционального диалога с потенциальным пилотным партнёром.

\textbf{Перспективы развития.} Основные направления дальнейшей работы включают: завершение разработки веб-сервиса администрирования и проведение пилотного внедрения в Омском государственном аграрном университете с переводом учебного заведения с парсингового на платформенный контур серверной части; расширение мобильных клиентов отдельными экранами шести типов справочных сущностей --- учебных групп, преподавателей, аудиторий, факультетов, кафедр и корпусов; расширение веб-сервиса администрирования инструментами визуализации нагрузки преподавателей и учебных групп, отдельным экраном занятости аудиторий и подсистемой отчётов с выгрузкой в форматы XLS и PDF; реализацию системных push-уведомлений о заменах и переносах занятий средствами APNs и FCM; внедрение автоматизированного тестирования (контрактные тесты сетевого слоя на мобильных клиентах, сквозные тесты на веб-административной части) и непрерывной интеграции; реализацию системных виджетов рабочего стола в Android-приложении и динамической темизации Material~You; локализацию интерфейса на дополнительные иностранные языки; юридическое оформление команды и регистрацию программного обеспечения в Федеральной службе по интеллектуальной собственности по итогам грантового конкурса <<Студенческий стартап>>.

Реализованная платформа представляет собой готовую основу для перехода к промышленному SaaS-продукту, ориентированному в первую очередь на учебные заведения среднего и малого размера, не располагающие собственной серверной инфраструктурой и квалифицированным ИТ-подразделением. Сочетание подтверждённой устойчивой эксплуатации, формирующейся институциональной апробации и проработанной дорожной карты дальнейшего развития позволяет рассматривать настоящую работу как самостоятельный технико-продуктовый результат, а не только как академическую разработку. Поставленная цель выпускной квалификационной работы достигнута.
